% Orden orientativo de ejecución; una fila por signo canónico.
\newcommand{\KalaTrazoFila}[2]{%
\KalaSimbolo[13mm]{#1}\quad\texttt{#1} &
\textbf{\csname KalaNombre#1\endcsname}
(\csname KalaTrayectos#1\endcsname\
\ifnum\csname KalaTrayectos#1\endcsname=1 trayecto\else trayectos\fi). #2\\[7pt]}
\begin{longtable}{@{}p{32mm}>{\raggedright\arraybackslash}p{\dimexpr\linewidth-32mm-2\tabcolsep\relax}@{}}
\toprule Modelo & Orden de trazado \\ \midrule\endhead
\bottomrule\endfoot
\KalaTrazoFila{K01}{1. Bajo por la izquierda y curvo hacia la derecha en la base.
2. Dibujo el ángulo superior desde la izquierda hacia la derecha y bajo.}
\KalaTrazoFila{K02}{1. Desde abajo a la izquierda, subo curvando por arriba hacia
la derecha y vuelvo a la izquierda hasta el centro. 2. Bajo desde ese extremo.}
\KalaTrazoFila{K03}{1. Trazo la curva ascendente de izquierda a derecha y termino
horizontal. 2. Dibujo la punta desde arriba a la izquierda, pasando por el
extremo derecho y bajando a la izquierda.}
\KalaTrazoFila{K04}{1. Desde abajo a la izquierda, dibujo una onda con arco
superior, valle y subida derecha. 2. Añado la punta superior de izquierda a derecha.}
\KalaTrazoFila{K05}{1. Bajo por la izquierda y giro a la derecha en la base.
2. Añado la barra superior hacia la derecha. 3. Añado la barra intermedia.}
\KalaTrazoFila{K06}{1. Arco izquierdo abierto, de arriba abajo.
2. Zigzag derecho: derecha, diagonal abajo a la izquierda y derecha.}
\KalaTrazoFila{K07}{1. Desde arriba a la izquierda, curvo por la derecha y vuelvo
a la izquierda por abajo. 2. Trazo la L interior bajando y girando a la derecha.}
\KalaTrazoFila{K08}{1. Cierro el rombo: izquierda, arriba, derecha, abajo,
izquierda. 2. Barra interior hacia la derecha hasta tocar el lado oblicuo,
sin sobresalir.}
\KalaTrazoFila{K09}{1. Corchete izquierdo: derecha, abajo, izquierda.
2. Barra derecha superior. 3. Barra derecha inferior.}
\KalaTrazoFila{K10}{1. Desde abajo a la izquierda, subo, redondeo la esquina,
sigo a la derecha, bajo y vuelvo hacia el interior a la izquierda.}
\KalaTrazoFila{K11}{1. Arco desde arriba a la izquierda por la derecha hasta
abajo. 2. Diagonal desde abajo a la izquierda hacia arriba a la derecha.}
\KalaTrazoFila{K12}{1. V desde arriba a la izquierda al centro inferior y arriba
a la derecha. 2. Bajo una cola desde el vértice inferior.}
\KalaTrazoFila{K13}{1. Desde arriba a la izquierda, voy al centro y vuelvo abajo
a la izquierda. 2. Del centro hacia la derecha y un retorno corto arriba a la izquierda.}
\KalaTrazoFila{K14}{1. U desde arriba a la izquierda, base curva y subida derecha.
2. Barra horizontal entre los lados. 3. Vertical central desde arriba, cruzando la barra.}
\KalaTrazoFila{K15}{1. La misma U. 2. Barra corta interior, separada de ambos
lados. Conservo esa separación y no añado la vertical de K14.}
\KalaTrazoFila{K16}{1. Desde abajo a la izquierda: subo, derecha, esquina
inclinada, bajo y cierro hacia la izquierda. 2. Barra interior hacia la derecha.}
\KalaTrazoFila{K17}{1. Ángulo izquierdo hacia el centro y de vuelta abajo a la
izquierda. 2. Ángulo derecho hacia el centro y de vuelta abajo a la derecha.
Dejo separados los vértices; no formo una X.}
\KalaTrazoFila{K18}{1. Desde arriba a la derecha, voy a la izquierda, bajo con
esquinas curvas y vuelvo a la derecha por abajo. 2. Ángulo interior hacia la
derecha y abajo a la izquierda.}
\KalaTrazoFila{K19}{1. Desde abajo a la izquierda, curvo por arriba hacia la
derecha, vuelvo en diagonal abajo a la izquierda y termino abajo a la derecha.}
\KalaTrazoFila{K20}{1. Barra superior hacia la derecha, diagonal abajo a la
izquierda y bajada vertical. 2. Barra intermedia hacia la derecha.}
\KalaTrazoFila{K21}{1. Escalera: subo, derecha, subo, derecha.
2. Punta de flecha desde arriba a la izquierda, extremo y abajo a la izquierda.}
\KalaTrazoFila{K22}{1. Desde arriba a la derecha, rodeo por arriba, izquierda y
abajo, subo por la derecha y termino con la barra hacia dentro. Todo es un
trayecto continuo; no añado una bajada ni repaso el arco.}
\end{longtable}
